%----------------------------------------------------------------------------------------
% Tailored CV — PhD in Super-Resolution Microscopy Methods
% Jonas Ries Lab, Max Perutz Labs, University of Vienna
%----------------------------------------------------------------------------------------

\documentclass[10pt,a4paper,sans]{moderncv}

\usepackage[utf8]{inputenc}
\usepackage[T1]{fontenc}

\moderncvstyle{classic}
\moderncvcolor{blue}

\usepackage[scale=0.78]{geometry}
\usepackage{ragged2e}

%----------------------------------------------------------------------------------------

\firstname{Kundai Farai}
\familyname{Sachikonye}

\address{Auf der Lichtung 19}{82178 Puchheim, Germany}
\mobile{(+49) 1626386344}
\email{kundai.sachikonye@bitspark.com}
\homepage{github.com/fullscreen-triangle}{github.com/fullscreen-triangle}

%----------------------------------------------------------------------------------------

\begin{document}

\makecvtitle

%----------------------------------------------------------------------------------------
%	EDUCATION
%----------------------------------------------------------------------------------------

\section{Education}

\cventry{2019--2021}{PhD candidate, Computational Lipidomics}{Technische Universität München}{}{}{
  Scientific computing pipelines, high-dimensional spectral data analysis,
  distributed ML. Departed to pursue independent AI and computational biology
  systems research.
}
\cventry{2018--2019}{MSc Computational Biology}{Universität Tübingen}{}{}{}
\cventry{2014--2017}{MSc Pharmaceutical Biotechnology}{HAW Hamburg}{}{}{}
\cventry{2011--2014}{BSc Biochemistry and Cell Biology}{Jacobs University Bremen}{}{}{}

%----------------------------------------------------------------------------------------
%	RESEARCH SOFTWARE
%----------------------------------------------------------------------------------------

\section{Research Software}

\cvitem{helicopter}{
  \textbf{Multi-Modal Microscopy Data Integration.}
  Combines fluorescence, brightfield, phase-contrast, and spectral imaging
  modalities through sequential constraint satisfaction. Each physical
  measurement is an independent exclusion operator; intersection of constraints
  reduces structural ambiguity from $\sim$10$^{60}$ configurations to unique
  solutions without dependence on single-modality resolution limits.
  Architecture for integrating SMLM localisation clouds with cryo-EM or
  confocal reference channels in structural cell biology.
  \href{https://github.com/fullscreen-triangle/helicopter}{github.com/fullscreen-triangle/helicopter}
}

\cvitem{verum}{
  \textbf{State Reconstruction from Sparse Observations (Geometric Model Fitting).}
  Trajectory completion operator $T = P \circ B \circ K$: Kirchhoff propagation
  at unobserved nodes, backward Viterbi MAP inference, thermodynamic feasibility
  projection onto convex set. Banach contraction ($\lambda\in[0.3,0.7]$) guarantees
  unique fixed-point convergence; full state recovered from 33\,\% surface
  observations, $<$10\,\% error, convergence to machine precision in 40 iterations.
  Validated on real telemetry (convergence residual $7.2\times10^{-9}$; fault
  detection 13 samples ahead of event; structural correlation $r=0.9997$).
  Direct analogue of LocMoFit geometric model fitting applied to SMLM localisation
  clouds.
  \href{https://github.com/fullscreen-triangle/verum}{github.com/fullscreen-triangle/verum}
}

\cvitem{purpose + homo-veloce}{
  \textbf{Surrogate Models / Simulator-Based Inference.}
  Backward Training Principle: train the surrogate at the hardest single-objective
  instance (pure-intent regime, full envelope, unambiguous feedback);
  every less-constrained sub-regime inherits $\epsilon$-competence by feasibility
  inclusion, not generalisation. Forward Training Collapse Theorem establishes a
  provably positive VC failure-gap for the conventional direction; Sample Complexity
  Theorem gives $(N+1)\times$ savings for $N$-layer cascades.
  homo-veloce distils Gaussian process + GNN physics solvers into lightweight
  1D CNNs, LSTMs, Transformers for real-time inference --- the same pattern as
  simulator-based inference for fluorophore localisation.
  \href{https://github.com/fullscreen-triangle/purpose}{github.com/fullscreen-triangle/purpose}
}

\cvitem{borgia}{
  \textbf{First-Principles Molecular Spectroscopy.}
  Derives electronic transitions, absorption, emission, and vibrational modes of
  molecules from bounded phase space geometry --- zero empirical parameters.
  Triple Equivalence Theorem (oscillatory dynamics $=$ categorical state
  enumeration $=$ partition operations) makes a GPU shader formally equivalent
  to a physical spectrometer. Validated against NIST: nine elements, six molecules,
  self-consistency errors $<$1.63\,\%.
  Applicable to rational fluorophore selection for multi-colour SMLM: predict
  excitation/emission spectra, extinction coefficients, and Stokes shifts
  from molecular structure without empirical measurement.
  Live: \href{https://honjo-masamune.vercel.app/tools}{honjo-masamune.vercel.app/tools}.
  \href{https://github.com/fullscreen-triangle/borgia}{github.com/fullscreen-triangle/borgia}
}

\cvitem{lagrangian}{
  \textbf{Optical Physics from First Principles.}
  Resonant Stack theorem proves Beer-Lambert, Kirchhoff circuit law, and
  thermal emission as formally identical operations ($T=1/e=0.3679$ in all
  three representations). First-principles model of light propagation through
  biological tissue applicable to PSF characterisation and aberration
  correction in thick-sample 3D SMLM. Validated: $G$ derived to machine
  epsilon (43/43 benchmarks); GPU shader telescope 10/10 ($6.9\times10^{-5}$
  median error).
  Live: \href{https://space-observatory-xi.vercel.app/instruments/resonant-stack}{space-observatory-xi.vercel.app}.
  \href{https://github.com/fullscreen-triangle/lagrangian}{github.com/fullscreen-triangle/lagrangian}
}

\cvitem{pdsvm}{
  \textbf{Scientific Database Search via Bounded Oscillatory Systems.}
  Purpose-Driven Shader VM: every stable data object is a bounded oscillatory system
  described by S-entropy coordinates $(S_k, S_t, S_e) \in [0,1]^3$;
  probe operator $\Pi_P(T)$ converges geometrically ($L\approx 0.89$) to unique
  stable slice $T^*$ (forcing conditions, not statistics).
  Three exhaustive GPU shader primitives (Completeness Theorem):
  SpectralDP (cosine similarity kernel), RayMarch (geodesic integral),
  TrieTraversal ($O(k)$ independent of $N$; $373\times$ at $N=10^5$).
  Direct isomorphism with LocMoFit: localization cloud = data; target model = purpose
  point $P$; stable slice = protein complex geometry satisfying structural constraints.
  10/10 validation; $<0.01\,\mathrm{ms}$ per iteration; WebGPU browser-deployable.
}

\cvitem{lavoisier}{
  \textbf{Spectral Data Analysis Pipeline.}
  Dual-pipeline MS data analysis: numerical pipeline + CNN visual pipeline
  converting spectra to temporal image sequences (PyTorch); 98.9\,\% feature
  extraction accuracy against NIST reference library; Ray distributed computing;
  $>$100\,GB dataset support. Signal processing architecture applicable to
  large-scale SMLM localisation file analysis.
  \href{https://github.com/fullscreen-triangle/lavoisier}{github.com/fullscreen-triangle/lavoisier}
}

%----------------------------------------------------------------------------------------
%	EXPERIENCE
%----------------------------------------------------------------------------------------

\section{Experience}

\cventry{2021--present}{Independent AI \& Computational Biology Developer}{Bitspark GmbH}{}{}{
  Full-time independent development of computational biology frameworks: multi-modal
  microscopy analysis, surrogate model training, spectral data pipelines, geometric
  model fitting. All systems publicly documented and validated.
}

\cventry{2019--2021}{Computational Lipidomics}{TUM / Bayerisches Forschungsinstitut}{Munich}{}{
  Mass spectrometry-based lipid species characterisation at scale: peak detection,
  ion annotation (LIPID MAPS, LipidBlast), multi-omics statistical modelling,
  federated learning across distributed clinical sites. Spectral data analysis of
  $>$100\,GB mzML datasets.
}

\cventry{2019}{Process Optimisation}{Fraunhofer Institut IGB}{Stuttgart}{}{
  Real-time sensor data acquisition and ML-based optimisation for industrial
  photobioreactors. Dynamic programming and LP methods. Python, C.
}

%----------------------------------------------------------------------------------------
%	TECHNICAL SKILLS
%----------------------------------------------------------------------------------------

\section{Technical Skills}

\cvitem{ML / Inference}{
  PyTorch, scikit-learn, XGBoost; 1D CNNs, LSTMs, Transformers, GNNs, Gaussian
  processes; simulator-based / likelihood-free inference; variational Bayes;
  knowledge distillation from physics-based solvers
}

\cvitem{Computational Physics}{
  Bounded phase space geometry; spectral analysis (DFT, Fourier optics);
  Beer-Lambert optical modelling; geometric model fitting; Banach contraction
  mappings; graph Laplacian methods; Kirchhoff circuit analogs
}

\cvitem{Data / HPC}{
  Ray, Dask, distributed computing; HDF5, mzML; memory-mapped I/O for
  large-scale spectral datasets; Docker, Kubernetes, FastAPI
}

\cvitem{Languages}{Python, Rust, TypeScript/Node.js, R, Java, C, SQL}

%----------------------------------------------------------------------------------------
%	LANGUAGES
%----------------------------------------------------------------------------------------

\section{Languages}

\cvitemwithcomment{English}{Native / Fluent}{}
\cvitemwithcomment{German}{Conversational}{Living in Germany since 2011}

%----------------------------------------------------------------------------------------

\renewcommand{\listitemsymbol}{-~}

\end{document}
